\chapter{Методологија}

Трудот ќе го следи развојот на Logche како проект од областа на софтверското инженерство.
Методологијата ќе биде итеративна: најпрво се дефинираат основните кориснички сценарија, потоа се развиваат прототипи за парсирање и чување податоци, а на крај се евалуираат точноста, употребливоста и ограничувањата на решението.

Развојниот процес ќе биде документиран во `CHANGES.md` и во подолги развојни белешки кога има значајни архитектурни одлуки или експериментални резултати.
Во главниот текст на трудот овие белешки ќе се користат како материјал за формално објаснување, а не како дневник.

Евалуацијата ќе се дефинира според зрелоста на системот.
За парсерот може да се користат множества од пример-записи со очекувани структурирани излези.
За локалните модели може да се мерат точност, латентност, потрошувачка на ресурси и квалитет на експанзија на кратките записи.
